\documentclass[10pt,twocolumn]{article}

\usepackage[letterpaper,margin=0.72in,columnsep=0.24in]{geometry}
\usepackage{times}
\usepackage{array}
\usepackage{balance}
\usepackage{booktabs}
\usepackage{caption}
\usepackage{enumitem}
\usepackage{fancyhdr}
\usepackage{float}
\usepackage{microtype}
\usepackage{seqsplit}
\usepackage{tabularx}
\usepackage{titlesec}
\usepackage{url}
\usepackage[hidelinks]{hyperref}

\hypersetup{
  pdftitle={Failure-Aware Agentic Recovery for OCR-RAG: Verified Experimental Status},
  pdfauthor={FAAR project}
}

\setlength{\headheight}{12pt}
\setlength{\parindent}{1em}
\setlength{\parskip}{0pt}
\setlength{\tabcolsep}{4pt}
\emergencystretch=1.5em
\setlist{nosep}
\captionsetup{font=small,labelfont=bf}
\newcolumntype{Y}{>{\raggedright\arraybackslash}X}
\newcolumntype{L}[1]{>{\raggedright\arraybackslash}p{#1}}
\newcommand{\hash}[1]{\texttt{\seqsplit{#1}}}

\titleformat{\section}{\large\bfseries}{\thesection.}{0.45em}{}
\titleformat{\subsection}{\normalsize\bfseries}{\thesubsection.}{0.35em}{}
\titlespacing*{\section}{0pt}{1.4ex plus .3ex minus .2ex}{.6ex}
\titlespacing*{\subsection}{0pt}{1.1ex plus .2ex minus .2ex}{.35ex}

\pagestyle{fancy}
\fancyhf{}
\fancyhead[L]{\small FAAR to AAAI}
\fancyhead[R]{\small Verified experimental status}
\fancyfoot[C]{\small\thepage}
\renewcommand{\headrulewidth}{0.3pt}
\renewcommand{\footrulewidth}{0pt}

\title{Failure-Aware Agentic Recovery for OCR-RAG:\\
Verified Experimental Status}
\author{FAAR project}
\date{8 August 2026}

\begin{document}

\maketitle

\begin{abstract}
OCR errors can make document retrieval and question answering fail even when the relevant page is present. FAAR is a selective recovery controller that diagnoses likely OCR failures and chooses among direct answering, retrieval backtracking, text correction, and visual fallback. This paper-style status report records the evidence available before the planned full benchmark runs. The implementation includes a fixed OHR-Bench split, pinned model revisions, complete-document retrieval controls, provenance logging, and an end-to-end smoke path. The current empirical quality numbers come from an offline 40-example mock evaluation, not from the planned paper baselines. All five planned baselines, B0 through B4, currently have zero paper-eligible result files. On the offline slice, every compared profile obtains EM $=0.1250$ and F1 $=0.2036$, with zero delta against the naive RAG profile. The available evidence supports selective routing and diagnosis-policy consistency, but not an answer-quality improvement claim. Full OHR OCR, ArXivQA page confirmation and OCR, MP-DocVQA acquisition, and the ordered baseline runs remain outstanding.
\end{abstract}

\textbf{Keywords:} OCR-RAG, document question answering, selective recovery, multimodal retrieval, experimental reproducibility

\section{Introduction}

Document question answering systems often treat OCR output as if it were a faithful representation of the page. That assumption fails on tables, formulas, multilingual text, and visually structured pages. A retrieval failure can therefore be caused by the retriever, the OCR text, or the answer extraction step. FAAR addresses this problem with a quality gate and a typed recovery policy.

The intended research claim is that selective recovery can improve answerability while using less multimodal computation than an always-visual pipeline. The evidence available at this stage is narrower. The implementation and experiment controls are substantially prepared, but the full paper protocol has not yet produced B0 to B4 results. This report makes that boundary explicit.

Two evaluations are separated throughout the paper. First, an earlier prototype evaluation compares five FAAR profiles on a 40-example offline mock slice. Second, the AAAI execution protocol defines real OHR-Bench, MP-DocVQA, and full-paper ArXivQA runs. Prototype profile numbers are reported as engineering evidence. They are not used as substitutes for the planned paper baselines.

\section{Experimental protocol}

\subsection{Datasets and preparation}

The OHR-Bench split is fixed with seed 42. It contains 5,948 training questions, 1,274 validation questions, and 1,276 test questions. The validation split covers 549 documents and 7,033 complete-document pages. The test split covers 567 documents and 6,842 pages. Retrieval corpora are required to contain the complete document rather than only the gold evidence pages.

The supervisor-approved ArXivQA preparation uses 500 questions across 497 full papers. All 8,955 pages were downloaded, parsed, rendered, and hash-validated. The inventory remains preparation-only: all 500 page mappings are unresolved and no pinned GOT-OCR page outputs have been registered. MP-DocVQA validation is not yet local; the required download is the authenticated RRC DocVQA Task 4 source.

The locked stack uses GPT-4o-2024-11-20 for the VLM, NV-Embed-v2 for embeddings, BGE reranker-v2-m3 for reranking, GOT-OCR 2.0 for OCR, Docling for PDF preparation, ColPali and VisRAG for visual retrieval, and ByT5 for correction. Exact revisions and source fingerprints are listed in Appendix~\ref{app:evidence}.

\subsection{Baseline sequence}

The planned baselines must run in the following order. Each result is saved before the next run begins.

\begin{table}[H]
\centering
\caption{Planned AAAI baseline sequence.}
\label{tab:baselines}
\small
\begin{tabularx}{\columnwidth}{@{}lY@{}}
\toprule
B0 & Text-only RAG; gate off; recovery off. \\
B1 & Always-VLM recovery. \\
B2 & Random recovery with diagnosis disabled. \\
B3 & ColPali visual retrieval. \\
B4 & VisRAG visual retrieval. \\
\bottomrule
\end{tabularx}
\end{table}

The gate threshold is tuned only after a real B0 validation result exists. The annotation study, main FAAR runs, ablations, confidence intervals, and paper tables depend on this sequence. B2 is a required comparison for the diagnosis-module claim and remains pending.

\section{Results available to date}

\subsection{Paper-protocol baseline status}

The repository contains no \path{results/b0.json}, \path{results/b1.json}, \path{results/b2.json}, \path{results/b3.json}, or \path{results/b4.json}. The only B0-labelled file is a one-document smoke payload. It uses one validation example, lexical mock retrieval, and no API request. The file explicitly marks itself as a smoke test and as not a paper result.

\begin{table}[t]
\centering
\caption{B0 smoke payload. This is not a paper baseline.}
\label{tab:smoke}
\small
\begin{tabular}{@{}lr@{}}
\toprule
Examples & 1 \\
EM & 0.0 \\
F1 & 0.0 \\
VLM rate & 0.0 \\
Harm rate & 0.0 \\
API requests & 0 \\
Cost (USD) & 0.0 \\
Runtime (s) & 0.0226 \\
\bottomrule
\end{tabular}
\end{table}

Thus, the paper-protocol baseline completion count is zero of five. The smoke result validates pipeline wiring and serialization. It does not estimate dataset performance.

\subsection{Offline prototype evaluation}

The current prototype artifacts contain a 40-example mock evaluation with five profiles. The same examples and metric definitions were used for every profile. The run used a mock visual backend with API calls disabled.

\begin{table}[H]
\centering
\caption{Offline mock profile evaluation on 40 examples per profile. These values are prototype evidence, not B0 to B4 paper results.}
\label{tab:profiles}
\scriptsize
\begin{tabularx}{\columnwidth}{@{}Yrrrrrr@{}}
\toprule
Profile & $n$ & NDCG@5 & Recall@5 & EM & F1 & Fallback \\
\midrule
FAAR full & 40 & 0.4789 & 0.5250 & 0.1250 & 0.2036 & 0.0750 \\
FAAR no backtrack & 40 & 0.4789 & 0.5250 & 0.1250 & 0.2036 & 0.0750 \\
FAAR no diagnosis & 40 & 0.4789 & 0.5250 & 0.1250 & 0.2036 & 0.0000 \\
FAAR no VLM & 40 & 0.4789 & 0.5250 & 0.1250 & 0.2036 & 0.0000 \\
Naive RAG & 40 & 0.4789 & 0.5250 & 0.1250 & 0.2036 & 0.0000 \\
\bottomrule
\end{tabularx}
\end{table}

All profile deltas against \texttt{naive\_rag} are zero for EM, F1, NDCG@5, and Recall@5. The profile comparison therefore does not show a quality advantage for typed recovery on this slice.

For \texttt{faar\_full}, 21 of 40 examples followed direct answering, 14 used text correction, 3 invoked visual fallback, and 2 retried retrieval. The non-pass rate was 0.475. The recorded diagnosis-policy consistency was 1.0.

\subsection{Claim boundary}

The current evidence supports two narrow statements. First, the mock profile uses visual fallback selectively, at a rate of 0.075, rather than treating every example as visual. The claim artifact estimates 0.925 savings relative to an always-on visual profile. Second, the diagnosis categories align with the executed routing policy on the same slice.

The evidence does not support an answer-quality gain. The typed-recovery delta against \texttt{naive\_rag} is 0.0 for both F1 and EM. The selective visual profile also has zero F1 and EM delta against the no-VLM profile. These findings should be presented as limitations of the current slice, not as evidence that the method improves accuracy.

\section{Readiness of the full experiment}

\begin{table}[H]
\centering
\caption{Execution status by planned phase. ``Implemented'' refers to code and contracts; it does not mean that the paper experiment is complete.}
\label{tab:phases}
\scriptsize
\begin{tabularx}{\columnwidth}{@{}lL{0.25\columnwidth}Y@{}}
\toprule
Phase & Status & Evidence and remaining condition \\
\midrule
0 Setup & Implementation complete; assets incomplete & Fixed split, loaders, provenance, asset contracts, and smoke path exist. Full OHR, MP-DocVQA, and ArXivQA assets are not ready. \\
1 Baselines & Pending & The B0 to B4 sequence is specified; 0 of 5 result files exists. \\
2 Gate & Code complete; validation pending & Threshold search is implemented, but a real B0 validation result is required. \\
3 Annotation & Tooling complete; study pending & Sampling, OCR extraction, Label Studio export, and kappa checks exist. Human labels do not. \\
4 Main FAAR & Pending assets and credentials & Runners and complete-asset guards exist; no full dataset result is present. \\
5 Ablations & Code complete; runs pending & A1 to A4 routing variants exist, but depend on the locked gate and benchmark assets. \\
6 Analysis & Code complete; inputs pending & Confidence intervals, failure analysis, harm rate, and Pareto export exist. Real result JSONs do not. \\
7 Paper tables & Inputs pending & Table and figure generation exists; no paper-eligible result tables have been generated. \\
\bottomrule
\end{tabularx}
\end{table}

\subsection{Preparation cost}

The one-page OHR smoke run used MPS. Docling took 8.91 seconds, rendering took 0.85 seconds, and GOT-OCR took 1,484.05 seconds per page. Peak RSS was 1.72 GiB, the PNG was 609,929 bytes, the OCR output was 7,100 bytes, and the model cache was approximately 2.12 GiB. A linear projection gives approximately 2,899 OCR-hours for OHR validation and 2,821 OCR-hours for OHR test. These figures explain why a GPU-backed batch plan is required; they are not completed dataset runtimes.

\subsection{ArXivQA readiness}

The 500 staged ArXivQA questions cover 497 papers and 8,955 rendered pages. Candidate matching accepted zero mappings and left all 500 unresolved. The paper inventory has zero OCR entries, state \texttt{preparation\_only}, and \texttt{paper\_ready=false}. The withdrawn paper \texttt{2304.04203} has a documented \texttt{v1} PDF override with a figure-match score of 0.7699745893478394; this is a provenance score for selecting the historical PDF, not an evaluation result.

\section{Discussion}

The implementation work is ahead of the empirical work. The controller, routing policies, asset contracts, and provenance checks are sufficiently concrete to support the planned evaluation. The available prototype results are more useful as a diagnostic than as a headline result: extraction changes lifted all profiles to the same EM and F1, while the recovery policies changed routing and fallback usage without changing quality.

This distinction also resolves the mixed phase language in the project documentation. The earlier phase reports describe prototype implementation milestones. The AAAI execution plan describes real benchmark runs. A statement that prototype Phases 0 to 4 are complete must not be read as a statement that the AAAI B0 to B4 baselines are complete.

There are three immediate threats to validity. The offline slice has only 40 examples per profile and uses mock visual execution. The full OCR assets are not yet available at practical throughput. Finally, ArXivQA is currently a preparation inventory with unresolved mappings and no page-level GOT-OCR provenance. None of these limitations can be repaired by changing the prose.

\section{Next experiment}

The next work should follow the fixed order below:

\begin{enumerate}[leftmargin=1.4em]
  \item Provision a GPU workflow for complete OHR validation and test OCR.
  \item Confirm all 500 ArXivQA page mappings and run pinned GOT-OCR with per-page provenance.
  \item Download and checksum MP-DocVQA validation from the RRC Task 4 source.
  \item Run B0, B1, B2, B3, and B4, saving each result before continuing.
  \item Tune and lock the validation gate after real B0 exists.
  \item Run annotation, FAAR, ablation, confidence-interval, failure-type, and cost analyses.
  \item Generate paper tables only from those saved result artifacts.
\end{enumerate}

\section{Conclusion}

FAAR is implemented as a selective, typed recovery controller for OCR-RAG. The current offline evidence shows selective routing and internally consistent diagnosis, but no answer-quality improvement over the naive profile. The full paper claim remains untested because B0 to B4 have not been executed on the locked benchmark protocol. The next milestone is therefore a reproducible GPU-backed baseline run, with B2 included, rather than another change to the claim language.

\balance
\clearpage
\onecolumn
\appendix

\section{Evidence map}
\label{app:evidence}

The main text reports only values that can be traced to the following artifacts. Paths are included here to keep the paper body readable.

\begin{table}[H]
\centering
\caption{Primary evidence artifacts.}
\small
\begin{tabularx}{\textwidth}{@{}L{0.40\textwidth}Y@{}}
\toprule
Artifact & Role in this report \\
\midrule
\path{PLAN.md} & AAAI phase status, implementation-versus-execution distinction, and blockers. \\
\path{docs/faar-aaai-plan.md} & Ordered B0 to B4 commands and completion checklist. \\
\path{RUNLOG.md} & Split counts, smoke measurements, asset preparation, verification, and blockers. \\
\path{results/smoke/b0_one_doc.json} & One-document B0 smoke payload; explicitly not a paper result. \\
\path{artifacts/phase3/metrics_summary.json} & Current offline 40-example profile metrics. \\
\path{artifacts/phase4/claim_assessment.json} & Supported and unsupported prototype claims. \\
\path{artifacts/phase4/profile_deltas_vs_baseline.csv} & Profile deltas against \texttt{naive\_rag}. \\
\path{config/model_revisions.json} & Exact revisions for the six locked Hugging Face models. \\
\path{config/arxivqa_source_lock.json} & ArXivQA source revisions, hashes, and the withdrawn-paper override. \\
\path{data/benchmark_prep/arxivqa/paper_inventory.json} & Rendered-page inventory and OCR readiness. \\
\path{data/benchmark_prep/arxivqa/evidence_candidates.json} & Confidence candidates; all 500 mappings remain unresolved. \\
\path{docs/reports/supervisor_progress_report.md} & Prototype mock command and 40-example result interpretation. \\
\bottomrule
\end{tabularx}
\end{table}

\section{Locked sources and revisions}

The official ArXivQA source is locked to revision \hash{85a6dca0e2bdc6f0268ae519be8913f83a83cafd} with SHA-256 \hash{d66cee08f152747ea58ce7851310c97229dcc6103901b3041ea988ad11eb1b96}. The ViDoRe source is locked to revision \hash{b8a106812c8682bab08935cf5d1b4566c82562de} with SHA-256 \hash{a45cb85bb49dce83b1beb851e312fe38179035151a55401a8801a2e53c1159d3}. The staged-row fingerprint is \hash{ccb60459021d211628c399682437841158e1ca238799a3cd81eef63026a7fa9a}.

The locked model revisions are NV-Embed-v2 \hash{3fa59658547db50a1e8e3346cf057fd0c77ed6ef}, BGE reranker \hash{953dc6f6f85a1b2dbfca4c34a2796e7dde08d41e}, GOT-OCR \hash{d3017ef2c2c1395888c8d635c5e0508bcb0ac78d}, ColPali \hash{56b5c732913abf3501ba208562df7b8f4edd5861}, VisRAG \hash{95ef596df871b606167cb7e4b7215caf1bfdf761}, and ByT5 \hash{68377bdc18a2ffec8a0533fef03b1c513a4dd49d}.

\section{Implementation verification}

The locked environment passed 136 tests with 4 warnings, and \texttt{pip check} reported no broken requirements. Python compilation passed and a local wheel build produced \texttt{faar-0.1.0-py3-none-any.whl}. The immutable split retained SHA-256 \hash{64583a532c5db5aa31e4cbb5cd9c7d894c7a2d5e8aa49f1a7f6041f54e714f53}. No coverage report or configured dead-code report is available. These checks establish implementation health; they do not establish benchmark completion.

\section{Recent implementation changes}

\begin{table}[H]
\centering
\caption{Recent committed changes supporting the experimental protocol.}
\small
\begin{tabular}{@{}llL{0.62\textwidth}@{}}
\toprule
Commit & Change & Experimental relevance \\
\midrule
\texttt{17128f0} & Ignore external inputs & Keeps downloaded benchmark data outside versioned source. \\
\texttt{5fffe18} & Lock ArXivQA sources & Records exact revisions and hashes. \\
\texttt{78ad6f2} & Prepare full-paper assets & Adds paper inventory and preparation support. \\
\texttt{fbab33e} & Add preparation CLI & Makes staging and preparation repeatable. \\
\texttt{2bd47e8} & Harden provenance & Fails closed on source and mapping inconsistencies. \\
\texttt{d4bc236} & Test asset preparation & Adds preparation and provenance coverage. \\
\texttt{2791339} & Enforce remap integrity & Protects complete-paper retrieval and mapping contracts. \\
\texttt{b69fb2b} & Record preparation evidence & Documents counts, readiness, and blockers. \\
\texttt{94e27a6} & Refresh phase status & Updates the execution handoff. \\
\texttt{8a7ef92} & Ignore build artifacts & Keeps generated output out of source control. \\
\bottomrule
\end{tabular}
\end{table}

\end{document}
